% Class-compatibility preamble retained for a reader excerpt of
%   Dave Witte Morris, "A Cohomological Proof that Real Representations of
%   Semisimple Lie Algebras Have Q-Forms", SIGMA 11 (2015) 034,
%   DOI 10.3842/SIGMA.2015.034, CC BY 4.0.
%
% The source pins the publisher class sigma, which is not part of the delivery
% baseline, so this file supplies the amsart baseline and then copies the author's
% own declaration block mechanically and unchanged:
%   - pinned-original-source.tex lines 11-31: the author's DeclareMathOperator
%     list, the equation numbering and the author's newtheorem declarations;
%   - pinned-original-source.tex lines 35-43, which the author writes immediately
%     after \begin{document} in the publisher layout: the author's own \Res,
%     \ResOne and \Lie macros and \allowdisplaybreaks.
% The publisher header commands of the sigma class (\FirstPageHeading,
% \ArticleName and the like) and the macro \PaperNumber they require are not
% reproduced, because this excerpt carries neither the journal front matter nor
% anything else that uses them.
% No mathematical macro, statement, symbol or assumption is introduced, changed
% or removed by this file.  The author's full DeclareMathOperator list of pinned
% source lines 11-19 is carried, including GL, SL, SO and SU: these four are not
% supplied by the delivery baseline, so leaving them out makes the author's own
% notation (for instance GL_n(Q) in the corollary on lattices) undefined.
\documentclass[10pt,a4paper]{amsart}
\usepackage{amssymb,amsmath,amsthm}
\usepackage{enumerate}
\usepackage{hyperref}

\DeclareMathOperator{\Cliff}{Clif\/f}
\DeclareMathOperator{\Gal}{Gal}
\DeclareMathOperator{\GL}{GL}
\DeclareMathOperator{\Inn}{Inn}
\DeclareMathOperator{\Mat}{Mat}
\DeclareMathOperator{\Out}{Out}
\DeclareMathOperator{\SL}{SL}
\DeclareMathOperator{\SO}{SO}
\DeclareMathOperator{\SU}{\mathbf{SU}}
\numberwithin{equation}{section}

\newtheorem{prop}{Proposition}[section]
\newtheorem{cor}[prop]{Corollary}
\newtheorem{lem}[prop]{Lemma}
{\theoremstyle{definition}
\newtheorem{defn}[prop]{Definition}
\newtheorem{notation}[prop]{Notation}
\newtheorem{rem}[prop]{Remark}
\newtheorem{ques}[prop]{Question}
\newtheorem{eg}[prop]{Example}
\newtheorem{erratum}[prop]{Erratum}}

\newcommand{\Res}[1]{\mathop{\mathbf{R}_{#1}}}
\newcommand{\ResOne}[1]{\mathop{\mathbf{R}^{(1)}_{#1}}}
\newcommand{\Lie}[1]{\mathfrak{\lowercase{#1}}}
\allowdisplaybreaks
\usepackage{cleveref}
\crefname{originallemma}{original lemma}{original lemmas}
\Crefname{originallemma}{Original lemma}{Original lemmas}
\makeatletter
\newlabel{RunivIff}{{2.1}{source}{}{sourcepointer-1}{}}
\newlabel{RunivIff@cref}{{[originallemma][0][]2.1}{[1][0][]source}{}{}{}}
\makeatother
\crefname{originaldefinition}{original definition}{original definitions}
\Crefname{originaldefinition}{Original definition}{Original definitions}
\makeatletter
\newlabel{RunivIff-D}{{4}{source}{}{sourcepointer-2}{}}
\newlabel{RunivIff-D@cref}{{[originaldefinition][0][]4}{[1][0][]source}{}{}{}}
\makeatother
\crefname{originaldefinition}{original definition}{original definitions}
\Crefname{originaldefinition}{Original definition}{Original definitions}
\makeatletter
\newlabel{RunivIff-DivAlg}{{3}{source}{}{sourcepointer-3}{}}
\newlabel{RunivIff-DivAlg@cref}{{[originaldefinition][0][]3}{[1][0][]source}{}{}{}}
\makeatother
\crefname{originallemma}{original lemma}{original lemmas}
\Crefname{originallemma}{Original lemma}{Original lemmas}
\makeatletter
\newlabel{SameFixed}{{2.3}{source}{}{sourcepointer-4}{}}
\newlabel{SameFixed@cref}{{[originallemma][0][]2.3}{[1][0][]source}{}{}{}}
\makeatother
\crefname{originalproposition}{original proposition}{original propositions}
\Crefname{originalproposition}{Original proposition}{Original propositions}
\makeatletter
\newlabel{so2k}{{5.1}{source}{}{sourcepointer-5}{}}
\newlabel{so2k@cref}{{[originalproposition][0][]5.1}{[1][0][]source}{}{}{}}
\makeatother
\title{R-universality of a unitary algebra over a quaternion division algebra in terms of the reduced norm and the even Clifford algebra}
\author{Dave Witte Morris}
\date{doi:10.3842/SIGMA.2015.034; CC BY 4.0}
\begin{document}
\maketitle
\noindent\textit{Editorial note.} All delivered mathematical text is the author's own, copied from the pinned publisher source sources/algebra-fl02-a-cohomological-proof-that-real-re/pinned-original-source.tex; no statement or proof step was rewritten, reordered or removed.\par\medskip
\noindent\textit{Editorial note.} Preamble provenance: the source pins the publisher class sigma, which is not part of the delivery baseline, so this card uses sources/algebra-fl02-a-cohomological-proof-that-real-re/retained-preamble.tex, which supplies the amsart baseline and then carries the author's own declarations from the pinned source (its header records the exact line ranges). The author's full DeclareMathOperator list is carried, including GL, SL, SO and SU, because the delivery baseline does not provide those four operator names. No macro was invented and no mathematical symbol was changed.\par\medskip
\noindent\textit{Editorial note.} Bibliography provenance: the citations of the counted proof are resolved with sources/algebra-fl02-a-cohomological-proof-that-real-re/bibliography.tex, an excerpt of the author's own thebibliography environment (pinned source lines 954-1016) copied character for character with all twelve entries kept; no entry was rewritten, reordered, merged, split or invented and no citation key was changed.\par\medskip
\noindent\textit{Editorial note.} External references are navigation pointers only; every deliberately kept step of the author's argument stays inside the counted spans.\par\medskip
\noindent\textit{Editorial note.} Selection provenance: this card was opened under the file-based control contract (CONTROL.md), with the mechanical pre-card gate run on its statement and proof interval before the card was written.\par\medskip
\noindent\textit{Editorial note.} The counted proof refers to the author's proposition labelled Beta=D; because that label contains an equals sign and the delivery format accepts only a restricted label alphabet, the proposition's statement is delivered verbatim as a context span instead of being declared as a navigation pointer.\par\medskip
\noindent\textit{Allowed standard prerequisites (editorial).}\begin{itemize}
\item The Tits algebra of the natural representation over a quaternion division algebra with the reversion anti-involution. Used inside the counted proof at source lines 741-742, attributed by the author to page 379 of the cited book KnusEtAl-BookOfInvols. Identifies the first Tits algebra with the quaternion division algebra itself.
\item The even Clifford algebra of a Hermitian form and the behaviour of the remaining Tits algebras according to the parity of the size. Used in the proof where the author compares the two remaining classes and uses that the centre of the spin group has exponent two in the even case (source lines 762-764). Produces the two alternatives of the statement, according to whether the size is even or odd.
\end{itemize}\medskip
\ifdefined\hypertarget\hypertarget{sourcepointer-1}{}\fi
\noindent\textit{Source reference.} Original lemma 2.1: sources/algebra-fl02-a-cohomological-proof-that-real-re/pinned-original-source.tex, source lines 416--416. Navigation pointer only: the counted proof invokes the author's equivalence for R-universality; that lemma has its own card in this lane and its content is declared as a named standard prerequisite.\par\medskip
\ifdefined\hypertarget\hypertarget{sourcepointer-2}{}\fi
\noindent\textit{Source reference.} Original definition 4: sources/algebra-fl02-a-cohomological-proof-that-real-re/pinned-original-source.tex, source lines 432--432. Navigation pointer only: the labelled item of that equivalence used by the counted proof.\par\medskip
\ifdefined\hypertarget\hypertarget{sourcepointer-3}{}\fi
\noindent\textit{Source reference.} Original definition 3: sources/algebra-fl02-a-cohomological-proof-that-real-re/pinned-original-source.tex, source lines 428--428. Navigation pointer only: the labelled item of that equivalence used by the counted proof.\par\medskip
\ifdefined\hypertarget\hypertarget{sourcepointer-4}{}\fi
\noindent\textit{Source reference.} Original lemma 2.3: sources/algebra-fl02-a-cohomological-proof-that-real-re/pinned-original-source.tex, source lines 471--471. Navigation pointer only: the author's lemma on the minimal inner extension; it has its own card in this lane and its content is declared as a named standard prerequisite.\par\medskip
\ifdefined\hypertarget\hypertarget{sourcepointer-5}{}\fi
\noindent\textit{Source reference.} Original proposition 5.1: sources/algebra-fl02-a-cohomological-proof-that-real-re/pinned-original-source.tex, source lines 674--674. Navigation pointer only: the counted proof refers to that proposition's setup; it is counted by this lane's card P1-fl02-so-even-clifford-split and its text is not reproduced here.\par\medskip
\begin{lem}
\label{SO(quat)}
Let $\Lie G = \Lie{SU}_k(B; D, \tau_r)$, where
\begin{itemize}\itemsep=0pt
\item $D$ is a~quaternion division algebra over~$\mathbb{Q}$,
\item $\tau_r$ is the reversion anti-involution, and
\item $B$~is a~$\tau_r$-Hermitian matrix in $\GL_k(D)$.
\end{itemize}
Then $\Lie G$ is $\mathbb{R}$-universal if and only if~$D$ does not split over~$\mathbb{R}$, and either
\renewcommand{\labelenumi}{{\rm (\theenumi)}}
\begin{enumerate}\itemsep=0pt
\item $k$ is even and the reduced norm of~$B$ $($calculated in the algebra $\Mat_{k \times k}(D))$ is
a~square in~$\mathbb{Q}$, or
\item $k$ is odd and $C^0_D(B, \tau_r)$ is split.
\end{enumerate}
\end{lem}

\begin{prop}[\protect{\cite[Corollary~3.5, \S~4.2, and Lemma~7.4]{Tits-RepLinIrred}}]\label{Beta=D}
Let~$L$ be the center of the algebra $D_{\Lie G, F}(\lambda):= \mathrm{End}_{\Lie G}({}^F \rho_\lambda)$ $($which,
by Schur's lemma, is a~division algebra$)$.
Then:
\renewcommand{\labelenumi}{{\rm (\theenumi)}}
\begin{enumerate}\itemsep=0pt
\item
\label{Beta=D-invt}
$\Gal(\overline{F}/L) = \{ \sigma \in \Gal(\overline{F}/F) \mid \sigma(\lambda) = \lambda \}$ $($for the $*$-action
of $\Gal(\overline{F}/F))$, and
\item $\beta_{\Lie G, L}(\lambda) = [D_{\Lie G, F}(\lambda)]$, after identifying $H^2(L; \boldsymbol\mu)$ with the
Brauer group of~$L$.
\end{enumerate}
\end{prop}

\begin{proof}
As in the proof of Proposition~\ref{so2k}, let $\lambda$, $\lambda_+$, $\lambda_-$ be dominant weights that represent the three nonzero
elements of~$Z^*$, with~$\lambda$ being the highest weight of the natural representation of~$\Lie G$ on
$(D\otimes_\mathbb{Q} \mathbb{C})^k$.

($\Rightarrow$) We have $D_{\Lie G, \mathbb{Q}}(\lambda) = D$~\cite[p.~379]{KnusEtAl-BookOfInvols}, so~$D$ does not
split over~$\mathbb{R}$ (see Lemma~\ref{RunivIff}(\ref{RunivIff-DivAlg})).
(Note that this implies $\Lie G \otimes_\mathbb{Q} \mathbb{R} \cong \Lie{SO}_k(\mathbb{H})$.)

Suppose~$k$ is even, so $\Lie G \otimes_\mathbb{Q} \mathbb{R}$ is inner.
Then $\Lie G$ must also be inner (see Lemma~\ref{SameFixed}).
Since~$k$ is even, this means that the reduced norm of~$B$ is a~square in~$\mathbb{Q}$.

Suppose~$k$ is odd.
Then $\Lie G \otimes_\mathbb{Q} \mathbb{R}$ is outer, so the weights $\lambda_+$ and~$\lambda_-$ are not f\/ixed by the
$*$-action.
Therefore $D_{\Lie G, \mathbb{R}}(\lambda_{\pm}) = \mathbb{C}$ (see Proposition~\ref{Beta=D}).
Since $D_{\Lie G, \mathbb{Q}}(\lambda_{\pm})$ is Brauer equivalent to $C^0_D(B,
\tau_r)$~\cite[p.~379]{KnusEtAl-BookOfInvols}, we conclude that $\Lie G$ is $\mathbb{R}$-universal if and only if this
algebra is split (see Lemma~\ref{RunivIff}(\ref{RunivIff-D})).

($\Leftarrow$) The proof for odd~$k$ is reversible, so let us assume~$k$ is even.
Then $D_{\Lie G, \mathbb{R}}(\lambda_{\pm}) = \mathbb{H}$ is nontrivial.
(For example, $D_{\Lie G, \mathbb{R}}(\lambda_{\pm})$ can be calculated by using~\cite[\S~5.5]{Tits-RepLinIrred}.)
Therefore $D_{\Lie G, \mathbb{Q}}(\lambda_{\pm})$ does not split over~$\mathbb{R}$.
Since $D_{\Lie G, \mathbb{Q}}(\lambda_{\pm})$ is a~quaternion algebra (because the center of $\mathbf{Spin}_{2k}$ has
exponent~$2$ when~$k$~is even), we conclude that $D_{\Lie G, \mathbb{Q}}(\lambda_{\pm}) \otimes_\mathbb{Q} \mathbb{R}
\cong D_{\Lie G, \mathbb{R}}(\lambda_{\pm})$, as desired.
\end{proof}

%% Bibliography excerpt retained for reader excerpts of
%   Dave Witte Morris, "A Cohomological Proof that Real Representations of Semisimple Lie Algebras Have Q-Forms",
%   SIGMA 11 (2015) 034, DOI 10.3842/SIGMA.2015.034.
%%
%% Every entry below is copied character for character from the author's own
% thebibliography environment of the pinned publisher source (lines 954-1016).
%% No entry was rewritten, reordered, merged, split or invented, and no citation
% key was changed.  All 12 entries are kept so that any citation of a delivered
%% excerpt resolves without guessing an entry.
\begin{thebibliography}{99}

\footnotesize\itemsep=0pt

\bibitem{DemazureGrothendieck-SGA3}
Demazure M., Grothendieck A., Sch\'emas en groups.~III. Structure des sch\'emas
  en groupes reductifs, S\'eminaire de g\'eom\'etrie alg\'ebrique du {B}ois
  {M}arie, 1962/1964, available at
  \url{http://library.msri.org/books/sga/sga/pdf/sga3-3.pdf}.

\bibitem{Harder-BerichtNeuere}
Harder G., Bericht \"uber neuere {R}esultate der {G}aloiskohomologie
  halbeinfacher {G}ruppen, \textit{Jber. Deutsch. Math.-Verein.} \textbf{70}
  (1968), 182--216, available at
  \url{http://www.digizeitschriften.de/en/dms/img/?PPN=PPN37721857X_0070&DMDID=dmdlog17}.

\bibitem{Kneser-LectGalCohoClassical}
Kneser M., Lectures on {G}alois cohomology of classical groups, \textit{Tata
  Institute of Fundamental Research Lectures on Mathematics}, Vol.~47, Tata
  Institute of Fundamental Research, Bombay, 1969, available at
  \url{http://www.math.tifr.res.in/~publ/ln/tifr47.pdf}.

\bibitem{KnusEtAl-BookOfInvols}
Knus M.A., Merkurjev A., Rost M., Tignol J.P., The book of involutions,
  \textit{American Mathematical Society Colloquium Publications}, Vol.~44,
  Amer. Math. Soc., Providence, RI, 1998.

\bibitem{Lam-IntroQuadForms}
Lam T.Y., Introduction to quadratic forms over f\/ields, \textit{Graduate Studies
  in Mathematics}, Vol.~67, Amer. Math. Soc., Providence, RI, 2005.

\bibitem{Morris-RRepsQForms}
Morris D., Real representations of semisimple {L}ie algebras have {${\mathbb
  Q}$}-forms, in Algebraic Groups and Arithmetic, Tata Institute of Fundamental
  Research, Mumbai, 2004, 469--490.

\bibitem{PlatonovRapinchukBook}
Platonov V., Rapinchuk A., Algebraic groups and number theory, \textit{Pure and
  Applied Mathematics}, Vol.~139, Academic Press, Inc., Boston, MA, 1994.

\bibitem{PrasadRapinchuk-Prescribed}
Prasad G., Rapinchuk A., On the existence of isotropic forms of semi-simple
  algebraic groups over number f\/ields with prescribed local behavior,
  \href{http://dx.doi.org/10.1016/j.aim.2006.01.001}{\textit{Adv. Math.}} \textbf{207} (2006), 646--660.

\bibitem{Raghunathan-ArithLattsSSGrps}
Raghunathan M.S., Arithmetic lattices in semisimple groups, \href{http://dx.doi.org/10.1007/BF02967980}{\textit{Proc.
  Indian Acad. Sci. Math. Sci.}} \textbf{91} (1982), 133--138.

\bibitem{Tits-Classification}
Tits J., Classif\/ication of algebraic semisimple groups, in Algebraic {G}roups
  and {D}iscontinuous {S}ubgroups ({P}roc. {S}ympos. {P}ure {M}ath., {B}oulder,
  {C}olo., 1965), Amer. Math. Soc., Providence, R.I., 1966, 33--62.

\bibitem{Tits-RepLinIrred}
Tits J., Repr\'esentations lin\'eaires irr\'eductibles d'un groupe r\'eductif
  sur un corps quelconque, \href{http://dx.doi.org/10.1515/crll.1971.247.196}{\textit{J.~Reine Angew. Math.}} \textbf{247} (1971),
  196--220.

\bibitem{Tits-StronglyInner}
Tits J., Strongly inner anisotropic forms of simple algebraic groups,
  \href{http://dx.doi.org/10.1016/0021-8693(90)90201-X}{\textit{J.~Algebra}} \textbf{131} (1990), 648--677.


\end{thebibliography}
\end{document}
